\clearpage
\VerdierPhysicalPageAnchor{0181}%
\begingroup\small\setlength{\parskip}{0pt}\setlength{\emergencystretch}{1.5em}
\setlength{\abovedisplayskip}{4pt plus 2pt minus 2pt}%
\setlength{\belowdisplayskip}{4pt plus 2pt minus 2pt}%

\VerdierSemanticAnchor{ch3d-1-2-10-proof-cont}
The complex $X(0,+\infty)$ (proof of (1.2.1)) is then quasi-isomorphic to
$X$, and the canonical morphism
$H^0(X(0,+\infty))\longrightarrow X(0,+\infty)$ is a quasi-isomorphism.
This shows that $X$ is isomorphic to $H^0(X)$ in $D(\mathcal A)$.

Let us prove b). The second equality is recalled for reference and follows
from the fact that $\mathcal A$ is a full subcategory of $D(\mathcal A)$.
A morphism from $X[-n]$ to $Y$ in $D(\mathcal A)$ has the form
$Q(m)Q(s)^{-1}$, where $m\colon Z^\bullet\longrightarrow Y$ is a morphism in
$K(\mathcal A)$ and $s\colon Z^\bullet\longrightarrow X[-n]$ is a
quasi-isomorphism in $K(\mathcal A)$. The morphism
$i\colon Z^\bullet(-\infty,n+1)\longrightarrow Z^\bullet$ (proof of
(1.2.1)) is then a quasi-isomorphism. Consequently,
\[
Q(m)Q(s)^{-1}=Q(m)Q(i)Q(i)^{-1}Q(s)^{-1}=Q(mi)Q(si)^{-1}.
\]
When $n$ is strictly negative, $mi=0$. Hence $Q(m)Q(s)^{-1}=0$.

\VerdierSemanticAnchor{ch3d-1-2-11}
\medskip\noindent\textbf{Definition 1.2.11.}\quad
We denote by $\mathcal A^{\operatorname{Ext}}$ the full subcategory of
$D(\mathcal A)$ defined by the translates of the objects of $\mathcal A$.

Thus the category $\mathcal A^{\operatorname{Ext}}$ is (unlike the
subcategory $\mathcal A$) a $\mathbb Z$-subcategory of $D(\mathcal A)$. An
object of $D(\mathcal A)$ is isomorphic to an object of
$\mathcal A^{\operatorname{Ext}}$ if and only if it has only one nonzero
cohomology object.

\VerdierSemanticAnchor{ch3d-1-3}
\subsection{The canonical \texorpdfstring{$\delta$}{delta}-functor.}

\VerdierSemanticAnchor{ch3d-1-3-1}
\noindent\textbf{1.3.1.}\quad
Let $\mathcal A$ be an abelian category and $\mathcal D$ a triangulated
category. A $\delta$-functor from $\mathcal A$ to $\mathcal D$ is an object
consisting of:

\noindent 1) A functor $G\colon\mathcal A\longrightarrow\mathcal D$.

\noindent 2) For every exact sequence
$S=0\longrightarrow X\longrightarrow Y\longrightarrow Z\longrightarrow0$
in $\mathcal A$, the choice of a morphism in $\mathcal D$:
\[
\delta(S)\colon G(Z)\longrightarrow G(X)[1].
\]
These data are subject to the following axioms:

\noindent a) For every morphism of exact sequences:

\VerdierSemanticAnchor{ch3d-1-3-1-exact-sequences}
\begin{center}
\resizebox{0.88\textwidth}{!}{%
$\begin{tikzcd}[ampersand replacement=\&,column sep=2.6em,row sep=2.2em]
S={}\&0\arrow[r]\&X\arrow[r]\arrow[d,"u"']\&Y\arrow[r]\arrow[d,"v"']\&Z\arrow[r]\arrow[d,"w"']\&0\\
S'={}\&0\arrow[r]\&X'\arrow[r]\&Y'\arrow[r]\&Z'\arrow[r]\&0
\end{tikzcd}$%
}
\end{center}
\endgroup
